\section{Introduction}
\label{sec:introduction}

A puzzle hunt is not a puzzle. This distinction, obvious to the puzzle hunt community,
carries real technical weight for anyone attempting to generate one with AI. A single puzzle
--- a crossword, a logic grid, a cryptic clue --- is a contained artifact with a well-defined
quality criterion: it either works mechanically or it does not. A puzzle hunt is a
\emph{multi-artifact creative system}: a set of thematically unified puzzles, each with a
confirmed final answer, feeding into a meta-puzzle whose extraction depends on the feeders'
specific answers, all wrapped in an editorial voice, a narrative world, and a delivery
package formatted for a specific medium (a website, a physical dossier, a print-and-play
booklet). The quality bar for a hunt is not just ``does each puzzle work?'' but ``does the
collection cohere? does the meta close? does the delivery hold together?''

This problem class --- multi-artifact creative system generation --- is not well addressed by
existing AI generation techniques. Single-shot prompting produces variable output
uncorrelated across artifacts; it has no mechanism for the kind of global constraint
propagation that puzzle hunt design requires. Agent-loop approaches
\cite{park2023generativeagents,nakajima2023babyagi} introduce iteration but do not
inherently impose the structured handoffs that keep downstream artifacts consistent with
upstream decisions. Procedural content generation (PCG) research has addressed related
problems in level generation \cite{shaker2016procedural,summerville2018procedural} and
narrative \cite{roemmele2015creative,kreminski2020why}, but these systems typically target
single-artifact or fixed-grammar outputs, not a heterogeneous pipeline of interdependent
creative artifacts.

The core difficulty is \emph{coordination across decisions}. A puzzle hunt requires that the
world-building done at Stage~1 constrain the puzzle types admitted at Stage~3, that the
puzzle answers selected at Stage~4 remain stable for the meta-puzzle designed at Stage~5,
and that the delivery format chosen at Stage~1 be honored by the assets produced at
Stage~9. Without explicit structure, AI generation at the scale of a hunt degenerates: the
world drifts, the meta fails to close, puzzles assume facts the world has not established,
and delivery assets reference puzzles that have been cut. Coherence collapse is the default
outcome of unconstrained multi-artifact generation.

\subsection*{Our Approach}

We address this problem with a structured, 11-stage AI-assisted pipeline in which each stage
produces a defined set of artifacts consumed by subsequent stages, and two stages function as
formal review gates with explicit pass criteria. The pipeline is not an agent loop and not a
single prompt chain. It is a decomposition of the full puzzle hunt production task into
stages that can be executed, reviewed, and retried independently, with structured artifacts
(markdown documents following defined templates) serving as the sole interface between
stages. Stage~6 (Authoring) in particular operates under dense human constraints --- a
fully-specified brief, a locked world, and a committed answer word --- so the pipeline is
better characterized as AI-assisted execution under specification than as open-ended AI
creativity. A \emph{world lock} mechanism --- a canonical world document treated as read-only
after Stage~1 --- prevents coherence drift across the pipeline by ensuring that all
downstream generation is grounded in a single approved source of truth
\cite{smith2010tanagra}.

Two stages are designated review gates:

\begin{itemize}
  \item \textbf{Stage~3 (Puzzle Pool Ranking):} An AI expert panel evaluates a pool of
  puzzle candidates against a published rubric and assigns numeric scores. Only candidates
  meeting the score threshold (or approved by human override) proceed to authoring. This
  gate is deliberately early: catching weak candidates before authoring investment is made
  is cheap; catching them after authoring is expensive.

  \item \textbf{Stage~7 (Editorial Review):} After authoring, each puzzle is submitted to a
  second AI panel for editorial review. The panel issues a pass/revise/cut verdict with
  specific, actionable feedback. Only puzzles passing Stage~7 proceed to integration. This
  gate prevents broken puzzles from entering the meta-puzzle and delivery package.
\end{itemize}

The pipeline is implemented as a library of 13 CLAUDE.md skill prompts --- persistent
markdown instruction files that configure an AI assistant to execute a specific pipeline
stage. The library runs without per-scenario code changes: the same skill prompts ran all
five evaluated scenarios, with only per-scenario world and scope documents varying.

\subsection*{Key Empirical Questions}

Given a structured pipeline with explicit gate conditions, three empirical questions become
tractable:

\begin{enumerate}
  \item \textbf{Does stage structure matter for output quality?} The 90.4\% first-pass pass
  rate across 52 puzzles provides an aggregate quality signal; failure distribution across
  stages provides the structural explanation.

  \item \textbf{Which stages are the highest-leverage intervention points?} If failures are
  uniformly distributed across stages, then gate placement is arbitrary. If failures
  concentrate at specific stages, that concentration identifies where pipeline investment has
  the highest return.

  \item \textbf{Is the pipeline transferable across themes and settings?} If pipeline quality holds only for
  the theme it was developed in, it is a one-off system, not a general contribution. If it
  transfers to new themes and settings without code changes, the pipeline specification itself is the
  contribution.
\end{enumerate}

We ground these questions in five fully-executed hunt scenarios spanning five distinct
thematic domains (a real-time strategy game, noir detective fiction, science fiction
role-playing, music, and print puzzles) and three delivery formats. The first scenario (Age
of Empires ``Wololo'') completed Stages~1 through~10 in 3~hours and 17~minutes (elapsed
wall-clock time; active human effort not separately measured) as measured by git commit
timestamps; subsequent scenarios ran in overlapping sessions, with the full
pipeline executed independently for each. Across all five scenarios, 52 puzzles were
authored, reviewed, and tested.

\subsection*{Contributions}

This paper makes the following contributions:

\begin{enumerate}
  \item \textbf{Pipeline specification.} We specify an 11-stage AI-assisted puzzle hunt
  pipeline with defined input/output artifacts per stage, explicit gate conditions at
  Stages~3 and~7, and an autonomy gradient characterizing the degree of AI versus human
  agency at each stage.

  \item \textbf{Cross-scenario empirical characterization.} Across five hunt scenarios and
  52 puzzles spanning five domains and three delivery formats, the pipeline achieves a
  90.4\% first-pass pass rate (47/52). We provide per-scenario and aggregate pass rate
  breakdowns. The first fully-executed scenario completed in 3h~17m (elapsed) from scope to ship-ready
  artifacts.

  \item \textbf{Failure analysis.} Five puzzles failed first-pass review. Three of the five
  failures occurred at Stage~8 (testing and integration), downstream of both review gates;
  zero failures were attributable to Stage~3 (pool gate). We categorize failure types
  (mechanism failure, answer coordination failure, construction error, delivery format gap)
  and map each to its earliest detectable stage.

  \item \textbf{Identification of highest-leverage stages.} The failure distribution
  confirms Stage~3 and Stage~7 as the pipeline's highest-leverage interventions. Stage~3
  gates the investment of authoring work: failing here is cheap. Stage~7 gates integration:
  failing here wastes authoring and blocks meta closure. The concentration of failures
  \emph{after} Stage~7 (at Stage~8), and zero failures at Stage~3, confirms that both gates
  are functioning as intended.

  \item \textbf{Open pipeline specification.} The full pipeline is specified in a public git
  repository: 142 commits over two calendar days, a skill library that grew from 0 to 13
  skills, and a reviewer profile library that grew from 12 to 29 profiles. The repository
  contains all scenario artifacts (scope documents, puzzle files, editorial review reports)
  sufficient to replicate or extend the pipeline to new domains without proprietary
  dependencies.
\end{enumerate}

\subsection*{Paper Organization}

Section~\ref{sec:related} situates the pipeline within prior work on PCG, AI-assisted
creative systems, and review-gate patterns in creative production. Section~\ref{sec:system}
describes the pipeline in full, including stage-by-stage artifact specifications, gate
conditions, and the implementation as a CLAUDE.md skill library. Section~\ref{sec:empirical}
presents the empirical characterization: scenario overview, pass rate analysis, stage timing,
failure analysis, and cross-theme transferability. Section~\ref{sec:discussion} discusses
the implications of the failure distribution, the role of the world lock mechanism, and
limitations of the current evaluation. Section~\ref{sec:conclusion} summarizes contributions
and identifies open questions.
